\begin{officialform}[BulletinEN]\begin{english}
\FormHeading
\FormLine{Project code: \ProjectCode}
\FormLine{Project title: \ProjectTitleEnglish}
\FormLine{Project period: \ProjectPeriodEnglish}
\FormLine{Total cost: 15,000,000 VND}
\FormLine{Project leader: \PrincipalName\ (\PrincipalPhone)}
\FormLine{Project members: Nguyễn Thanh Trọng, Cao Tường Hưng, Nguyễn Như Quỳnh, Ngô Hưng Thịnh.}
\FormLine{Advisor: \SupervisorName\ (\SupervisorEmail)}

\FormLine{Necessity of the project: Duplicate, unlocated or delayed rescue reports hinder coordination under intermittent connectivity.}

\FormLine{Objectives: Develop a prototype combining edge image inference, compact transmission and map-based event clustering.}

\FormLine{Methodology: Train MobileNetV3-Large and evaluate image classification, clustering, ranking and dispatch separately using saved artifacts and simulation datasets.}

\FormLine{Project activities: Image AI, structured urgency, persistent queues, adaptive sending, synchronization and dashboard development; conditional edge and copy-invariance audits.}

\FormLine{Research results: The original seed-1024 model achieves \NewImageAccuracy\% accuracy and \NewImageMacroFOne\% macro-F1 on 256 test images in CPU evaluation, with three critical errors. Runtime conversion, quantization and pruning were evaluated separately; quantization reduced model size but also accuracy. Phone benchmarks cover 32 configurations on four devices. A Flutter app, FastAPI/SQLite backend and dashboard were developed.}

\FormLine{Research new finding: RQ1--RQ3 reveal conditional safeguards and limitations without establishing general superiority or field rescue benefits.}

\FormLine{Application potentials: Research and workflow demonstration. Urgency integration, mobile provenance, repeated device/network/SMS tests and expert validation remain necessary.}
\end{english}\end{officialform}
